% Capítulo 7 — Resultados esperados. Criado em 2026-11-16: reúne a regra de
% decisão (antes seção da Metodologia) e os resultados esperados (antes no
% capítulo do cronograma). Texto e quadro são os da v14.

\chapter{Resultados esperados}
\label{cap:resultados}

\label{sec:regra}Sob a hipótese principal, espera-se $\beta > 0$
(Capítulo~\ref{cap:objetivos}). O que o ensaio entrega, porém, depende do que
a medição dos voos e o cálculo de poder mostrarem, e o Quadro~\ref{qua:regra}
fixa cada caso; o Apêndice~\ref{ap:regra} traz definições, critérios numéricos
e o enquadramento dos casos intermediários, entre eles os registros parciais
e os tardios. Os critérios foram escritos depois das rodadas exploratórias,
com os resultados delas à vista; nenhum é calculado com o desfecho do pós-lei. O efeito plausível é uma faixa: a
magnitude individual do análogo equatoriano, em três cenários que levam à
regra as ressalvas do Capítulo~\ref{cap:referencial}, multiplicada pela fração
de nascimentos exposta, em grade fixada no apêndice; decide a célula de 80
gramas e fração de 0,25.
Poder suficiente é condição necessária para o teste, não garantia de que o
efeito exista nessa magnitude. A data-limite para os registros é 31 de maio de
2027, fim do terceiro bimestre.

\begin{quadro}[htb]
\caption{\label{qua:regra}Regra de decisão: o que o ensaio entrega em cada cenário}
\centering
\footnotesize
\begin{tabular}{@{}p{0.3cm}p{4.4cm}p{2.6cm}p{7.3cm}@{}}
\toprule
 & Medição dos voos & Poder & O que o ensaio entrega \\
\midrule
1 & Voos documentados por município e mês, com pico sazonal & Detectável $\leq$ plausível &
Desenho principal, confirmatório, com municípios e estação medidos. \\
2 & Idem & Detectável $>$ plausível &
Medida dos voos; desenho principal reportado como limite, sem teste. \\
3 & Voos documentados, sem pico sazonal & Qualquer &
Medida dos voos; o calendário não se aplica; forma reduzida descritiva. \\
4 & Série completa sem voos no período de referência & Qualquer &
Medida do alcance da lei; forma reduzida descritiva; sem estimativa de efeito da pulverização. \\
5 & Voos na referência e nenhum em 2018 (série completa): a prática cessou antes da lei & Qualquer &
Medida do alcance da lei; $\beta$ só descritivo, pois não distingue a lei do fim anterior da prática; sem estimativa de efeito da pulverização. \\
6 & Registros não obtidos, ou com menos de dois anos completos, até a data-limite & Com a classificação provisória &
Desenho principal com os municípios com aeronave no Censo de 2006 e a quadra chuvosa, exploratório (o parâmetro do Capítulo~\ref{cap:metodologia} não é o alvo); cálculo de poder; medida parcial dos voos, se houver. \\
\bottomrule
\end{tabular}
\fonte{Elaboração própria. Definições e critérios no Apêndice~\ref{ap:regra}. A linha 4 prevalece sobre a 5, e a 5 sobre as linhas 1 a 3.}
\end{quadro}

O Apêndice~\ref{ap:regra} diz, linha a linha, o produto mínimo e por que ele basta, ou não, como ensaio. Em
qualquer linha, entregam-se (i) a descrição do hiato sazonal de peso ao nascer
antes da lei e (ii) o efeito mínimo detectável de cada análise; nas linhas 1 a
5, (iii) a medida municipal e mensal da aplicação aérea antes e depois da
proibição; só na linha 1, (iv) a estimativa do efeito, confrontada com as
análises complementares. Nas linhas 3 e 6 o ensaio não basta sozinho, e na 6
a medida dos voos é parcial ou falta.
